\documentclass[preview]{standalone}

\usepackage{amsmath}
\usepackage{amssymb}
\usepackage{stellar}
\usepackage{enumitem}
\usepackage{definitions}

\begin{document}

\id{linearalgebra-exercises-batch-8}
\genpage

\section{Exercises}

\begin{snippetexercise}{linear-algebra-batch-8-ex-1}{}
    Let \(k \in \realnumbers\) and consider the endomorphism \(f_k \colon \realnumbers^3 \fromto \realnumbers^4\)
    given by \[f_k(x,y,z) \triangleq (3x + 2y + z, x -y -z, 4x + ky - z)\]
    Determine for which \(k\), \(f_k\) is invertible.
\end{snippetexercise}

\begin{snippetsolution}{linear-algebra-batch-8-ex-1-sol}{}
    \todo
\end{snippetsolution}

\begin{snippetexercise}{linear-algebra-batch-8-ex-2}{}
    If we consider \(\complexnumbers\) as an \(\realnumbers\)-vector space, the map \(c\colon \complexnumbers \fromto \complexnumbers\), \(c(z) = \bar{z}\) is linear.
    \begin{enumerate}[label=\Alph*]
        \item Prove that \(c\) is invertible and that \(c^{-1} = c\).
        \item Geometrically describe the effect of \(c\) on the plane \(\complexnumbers = \realnumbers^2\).
        \item Find the matrix associated with \(c\) with respect to the basis \(\{1, i\}\) of \(\complexnumbers\) as an \(\realnumbers\)-vector space.
    \end{enumerate}
\end{snippetexercise}

\begin{snippetsolution}{linear-algebra-batch-8-ex-2-sol}{}
    \todo
\end{snippetsolution}

\begin{snippetexercise}{linear-algebra-batch-8-ex-3}{}
    Given a fixed \(m \in \realnumbers\), consider the linear map \(f_m\colon \realnumbers^2 \fromto \realnumbers^2\), \(f_m(x, y) = (x + my, y)\).
    \begin{enumerate}[label=\Alph*]
        \item Geometrically describe the effect of \(f_m\) on the plane \(\realnumbers^2\).
        \item Find the matrix \(A\) associated with \(f\) with respect to the standard basis of \(\realnumbers^2\).
        \item Prove that \(A\) is invertible and calculate \(A^{-1}\).
    \end{enumerate}
\end{snippetexercise}

\begin{snippetsolution}{linear-algebra-batch-8-ex-3-sol}{}
    \todo
\end{snippetsolution}

\begin{snippetexercise}{linear-algebra-batch-8-ex-4}{}
    Consider the linear map \(\Phi \colon \mathcal{P}_n(\realnumbers) \fromto \realnumbers^{n+1}\), \[\Phi(p) = \left(p(0), p'(0), \frac{p''(0)}{2!}, \dots, \frac{p^{(n)}(0)}{n!}\right),\] where \(p^{(k)}\) is the \(k\)-th derivative of \(p\), defined recursively by \(p^{(0)} = p\) and \(p^{(k)} = (p^{(k-1)})', k \geq 1\).
    \begin{enumerate}[label=\Alph*]
        \item Find the matrix associated with \(\Phi\) with respect to the bases \(\{1, x, \dots, x^n\}\) of \(\mathcal{P}_n(\realnumbers)\) and the standard basis of \(\realnumbers^{n+1}\).
        \item Prove that \(\Phi\) is invertible and explicitly write \(\Psi = \Phi^{-1} \colon \realnumbers^{n+1} \fromto \mathcal{P}_n(\realnumbers)\).
        \item Prove Taylor's formula: given \(p \in \mathcal{P}_n(\realnumbers)\), \[p(x) = p(0) + p'(0)x + \frac{p''(0)}{2!}x^2 + \dots + \frac{p^{(n)}(0)}{n!}x^n.\]
    \end{enumerate}
\end{snippetexercise}

\begin{snippetsolution}{linear-algebra-batch-8-ex-4-sol}{}
    \todo
\end{snippetsolution}

\begin{snippetexercise}{linear-algebra-batch-8-ex-5}{}
    Consider the linear map
    \(f \in \operatorname{End}(\mathcal{P}_2(\realnumbers))\) given by \(f(p(x)) = xp'(x)\).
    \begin{enumerate}[label=\Alph*]
        \item Find the matrix associated with \(f\) with respect to the basis \(\{1, x, x^2\}\) of \(\mathcal{P}_2(\realnumbers)\).
        \item Determine if there exists a basis \(\mathcal{B}\) of \(\mathcal{P}_2(\realnumbers)\) such that \(\mathcal{M}(f, \mathcal{B}) = I_3\).
    \end{enumerate}
\end{snippetexercise}

\begin{snippetsolution}{linear-algebra-batch-8-ex-5-sol}{}
    \todo
\end{snippetsolution}

\begin{snippetexercise}{linear-algebra-batch-8-ex-6}{}
    Consider the map \(\Psi \colon \mathcal{P}_3(\mathbb{F}) \fromto \mathcal{P}_3(\mathbb{F})\), \((\Psi(p))(x) = p(x + 1) - p(x)\). As an example, if \(p(x) = x^2 - 5x + 7\), \[(\Psi(p))(x) = p(x + 1) - p(x) = (x + 1)^2 - 5(x + 1) + 7 - (x^2 - 5x + 7) = 2x - 4.\]
    \begin{enumerate}[label=\Alph*]
        \item Prove that \(\Psi\) is linear.
        \item Determine the matrix associated with \(\Psi\) with respect to the basis \(\{1, x, x^2, x^3\}\) of \(\mathcal{P}_3(\mathbb{F})\).
        \item Determine whether \(\Psi\) is invertible.
    \end{enumerate}
\end{snippetexercise}

\begin{snippetsolution}{linear-algebra-batch-8-ex-6-sol}{}
    \todo
\end{snippetsolution}

\end{document}